\chapter{Fundamental Identities}
\label{ch:ch09-identities}

\section{A Small Toolkit from the Unit Circle}

The previous chapter discovered that the sum and difference formulas are
not arbitrary rules to be memorized, but consequences of a single
geometric fact: rotations compose by adding angles. This raises a broader
question worth pursuing before the book moves any further, namely which
other relationships among the trigonometric functions are similarly forced
by geometry rather than imposed by convention. The answer is that a
surprising share of the identities used throughout trigonometry follow from
only a few underlying ideas already established in this book: the equation
of the unit circle, the reciprocal relationships among the six
trigonometric functions, the quotient definition of tangent, the symmetry
properties of sine and cosine proved in Chapter~6, and the angle-addition
formulas derived in Chapter~8. The purpose of this chapter is to gather
these relationships into a compact toolkit that will support everything
that follows.

\section{The Pythagorean Identity}

The most important identity in the toolkit comes directly from the
definition of the unit circle itself. Recall from Chapter~4 that every
point on the unit circle satisfies $x^2+y^2=1$, while $x = \cos\theta$ and
$y = \sin\theta$.

\begin{theorem}[Pythagorean Identity]
$\sin^2\theta + \cos^2\theta = 1$ for every angle $\theta$.
\end{theorem}

\begin{proof}
Substituting $x=\cos\theta$ and $y=\sin\theta$ into $x^2+y^2=1$ gives
$\cos^2\theta + \sin^2\theta = 1$, which is the same statement with its two
terms reordered.
\end{proof}

It is worth pausing to notice what has, and has not, just been
established. This is not really a new fact about trigonometry; it is
simply the equation of the unit circle, restated in trigonometric notation.
No new geometry has been introduced here, only a familiar piece of
geometry seen through new symbols.

\section{Derived Pythagorean Identities}

The remaining two Pythagorean identities follow from the first by ordinary
algebra rather than by any further appeal to geometry.

\begin{theorem}
$1 + \tan^2\theta = \sec^2\theta$ whenever $\cos\theta \neq 0$.
\end{theorem}

\begin{proof}
Divide every term of $\sin^2\theta + \cos^2\theta = 1$ by $\cos^2\theta$:
\[
\frac{\sin^2\theta}{\cos^2\theta} + \frac{\cos^2\theta}{\cos^2\theta} =
\frac{1}{\cos^2\theta}.
\]
Recognizing $\tan\theta = \sin\theta/\cos\theta$ and $\sec\theta =
1/\cos\theta$ gives $\tan^2\theta + 1 = \sec^2\theta$.
\end{proof}

\begin{theorem}
$1 + \cot^2\theta = \csc^2\theta$ whenever $\sin\theta \neq 0$.
\end{theorem}

\begin{proof}
Divide every term of $\sin^2\theta + \cos^2\theta = 1$ by $\sin^2\theta$:
\[
1 + \frac{\cos^2\theta}{\sin^2\theta} = \frac{1}{\sin^2\theta}.
\]
Recognizing $\cot\theta = \cos\theta/\sin\theta$ and $\csc\theta =
1/\sin\theta$ gives $1 + \cot^2\theta = \csc^2\theta$.
\end{proof}

Neither of these two identities is an independent fact requiring separate
memorization; each is the original Pythagorean identity, divided through by
a different square.

\section{Reciprocal and Quotient Identities}

Several further identities restate definitions already introduced earlier
in the book, gathered here for reference alongside their more substantial
neighbors. The reciprocal identities record that
\[
\csc\theta = \frac{1}{\sin\theta}, \qquad
\sec\theta = \frac{1}{\cos\theta}, \qquad
\cot\theta = \frac{1}{\tan\theta},
\]
and the quotient identities record that
\[
\tan\theta = \frac{\sin\theta}{\cos\theta}, \qquad
\cot\theta = \frac{\cos\theta}{\sin\theta}.
\]
Each of these five equations is simply a definition from earlier chapters
written in the form of an identity, and none requires its own proof.

\section{Cofunction Identities}

The difference formula of Chapter~8 immediately produces another useful
family of identities.

\begin{theorem}
$\sin\left(\dfrac{\pi}{2} - \theta\right) = \cos\theta$.
\end{theorem}

\begin{proof}
Apply the difference formula $\sin(A-B) = \sin A\cos B - \cos A\sin B$ with
$A = \pi/2$ and $B = \theta$:
\[
\sin\left(\frac{\pi}{2}-\theta\right) = \sin\frac{\pi}{2}\cos\theta -
\cos\frac{\pi}{2}\sin\theta.
\]
Since $\sin(\pi/2) = 1$ and $\cos(\pi/2) = 0$, this reduces to
$\sin\left(\pi/2 - \theta\right) = \cos\theta$.
\end{proof}

The same substitution into the difference formula for cosine gives
$\cos\left(\pi/2 - \theta\right) = \sin\theta$, and dividing the two
results gives $\tan\left(\pi/2 - \theta\right) = \cot\theta$, with
corresponding identities holding for secant and cosecant. These are called
\emph{cofunction identities} because each function is transformed into its
complementary partner, and unlike the informal complementary-angle
observations sometimes made about right triangles, each one here follows
directly from the rotation-composition machinery of Chapter~8.

\section{Even and Odd Identities}

Chapter~6 established the symmetry properties of sine, cosine, and, by
extension, tangent; those results are now incorporated into the toolkit of
this chapter rather than re-derived. Cosine is even, $\cos(-\theta) =
\cos\theta$, while sine and tangent are odd, $\sin(-\theta) = -\sin\theta$
and $\tan(-\theta) = -\tan\theta$. These identities are frequently useful
when simplifying expressions that involve a negative angle, since they
allow the negative sign to be moved outside the function entirely.

\section{How to Prove a Trigonometric Identity}

Unlike an ordinary equation, an identity is established not by solving for
an unknown but by a chain of logically valid rewritings connecting one side
to the other, and a systematic approach makes this process considerably
easier to carry out reliably.

\begin{algorithmproc}{How to Prove a Trigonometric Identity}
Begin with whichever side of the proposed identity looks more complicated,
and rewrite any tangent, cotangent, secant, or cosecant appearing in it in
terms of sine and cosine whenever doing so seems useful. Apply the
Pythagorean identities to replace expressions such as $1 - \cos^2\theta$
with $\sin^2\theta$, or the reverse, whenever such a substitution
simplifies the expression. Continue simplifying algebraically, working on
one side of the proposed identity at a time rather than manipulating both
sides simultaneously, and never cross-multiply across an equation that has
not yet been established as true. Continue this process until the
expression you started from matches the other side exactly.
\end{algorithmproc}

\begin{example}
Prove that $\dfrac{1-\cos\theta}{\sin\theta} = \dfrac{\sin\theta}{1+\cos\theta}$.
Beginning with the left-hand side and multiplying numerator and denominator
by $1+\cos\theta$,
\[
\frac{1-\cos\theta}{\sin\theta} =
\frac{(1-\cos\theta)(1+\cos\theta)}{\sin\theta(1+\cos\theta)} =
\frac{1-\cos^2\theta}{\sin\theta(1+\cos\theta)}.
\]
By the Pythagorean identity, $1-\cos^2\theta = \sin^2\theta$, so
\[
\frac{1-\cos\theta}{\sin\theta} =
\frac{\sin^2\theta}{\sin\theta(1+\cos\theta)} = \frac{\sin\theta}{1+\cos\theta},
\]
which matches the right-hand side.
\end{example}

\begin{practicenow}
Prove that $\dfrac{\cos\theta}{1-\sin\theta} = \dfrac{1+\sin\theta}{\cos\theta}$,
by multiplying numerator and denominator of the left-hand side by
$1+\sin\theta$.
\end{practicenow}

\begin{sidebar}{Chinese Syllable Structure and Admissible Space}
Standard Chinese phonology offers a sharper illustration of a constrained
state space than a general appeal to phonotactics. A syllable can be
described, to a first approximation, as a tuple of five independent slots:
an initial consonant, a medial glide, a vowel nucleus, a final consonant or
glide, and a lexical tone. If each of the five slots could take any of its
individually possible values without restriction, the total number of
syllables would be the product of the five slot sizes, $n_1 n_2 n_3 n_4
n_5$. The actual inventory of syllables in the language is very much
smaller than this product, because entire combinations of slot values are
disallowed: certain initials never combine with certain medials, and
several finals occur with only a restricted subset of nuclei. The
language therefore occupies a genuinely smaller \emph{admissible space}
carved out of the larger \emph{theoretical space} that the unconstrained
product would predict, in exactly the sense introduced in Chapter~0's
discrete state spaces. The Pythagorean identity of this chapter performs
the same carving operation in a continuous setting: the theoretical space
of all pairs $(x,y) \in \mathbb{R}^2$ is far larger than the admissible
space of pairs that can actually occur as $(\cos\theta, \sin\theta)$,
which is restricted to the unit circle by the single constraint $x^2 +
y^2 = 1$. In both the syllable inventory and the trigonometric identity,
the interesting mathematical object is not the raw product of independent
possibilities but the much smaller admissible subset that a constraint
actually permits.
\end{sidebar}

\section{Identity Versus Equation}

A distinction is worth making explicit before proceeding further. An
\emph{identity} is a statement true for every value of the variable for
which both sides are defined; $\sin^2\theta + \cos^2\theta = 1$, for
instance, holds for every angle $\theta$ without exception. An
\emph{equation}, by contrast, is true only for specific values of the
variable; $\sin\theta = 1/2$ holds only for certain angles, and finding
those angles is a solving problem rather than a proving problem. Confusing
these two tasks --- attempting to "solve" an identity, or attempting to
"prove" an equation true for all $\theta$ when it plainly is not --- is one
of the most common sources of difficulty for students first encountering
this material, and the distinction will matter again directly when
Chapter~11 turns to solving trigonometric equations.

\section{Applications}

Trigonometric identities appear throughout mathematics, science, and
engineering: they simplify formulas in physics, support the analysis of
oscillatory systems, underlie techniques in signal processing, facilitate
coordinate transformations, and make possible the trigonometric
substitution techniques used later in calculus. The value of an identity
lies not merely in the fact that it is true, but in what that truth makes
possible: it allows one expression to be replaced by another that may be
easier to compute, simplify, or interpret, without changing the underlying
quantity being described.

\section{Looking Ahead}

At this point the book has assembled a substantial identity toolkit:
reciprocal identities, quotient identities, the three Pythagorean
identities, the cofunction identities, the even and odd identities, and the
sum and difference formulas of Chapter~8. Remarkably, this modest
collection already suffices to generate most of the remaining identities
encountered in elementary trigonometry. The next chapter shows how the
double-angle, half-angle, and product-to-sum formulas all arise naturally
from the tools developed here, so that what can appear, from a distance, to
be an enormous and disconnected collection of formulas turns out to be
different expressions of a comparatively small set of underlying
principles.

\section{Exercises}
\begin{exercise}
Starting from $\sin^2\theta + \cos^2\theta = 1$, derive $1 + \tan^2\theta =
\sec^2\theta$.
\end{exercise}

\begin{exercise}
Starting from $\sin^2\theta + \cos^2\theta = 1$, derive $1 + \cot^2\theta =
\csc^2\theta$.
\end{exercise}

\begin{exercise}
Prove that $\dfrac{\sin\theta}{\cos\theta} = \tan\theta$ using the quotient
identity.
\end{exercise}

\begin{exercise}
Prove that $\dfrac{1-\sin^2\theta}{\cos\theta} = \cos\theta$.
\end{exercise}

\begin{exercise}
Prove that $\dfrac{\cos^2\theta}{1-\sin\theta} = 1+\sin\theta$.
\end{exercise}

\begin{exercise}
Use the cofunction identities to evaluate $\sin\left(\dfrac{\pi}{2} -
\dfrac{\pi}{6}\right)$.
\end{exercise}

\begin{exercise}
Use the even and odd identities to simplify $\sin(-\theta) + \cos(-\theta)$.
\end{exercise}

\begin{exercise}
Determine whether $\cos\theta = 0$ is an identity or an equation, and
explain your reasoning.
\end{exercise}

\begin{exercise}
Explain why $(\sin\theta, \cos\theta) = (2,3)$ can never occur for any
angle $\theta$.
\end{exercise}

\begin{exercise}
Why is $\sin^2\theta + \cos^2\theta = 1$ better understood as a geometric
constraint than as an isolated algebraic formula?
\end{exercise}
